Segundo \cite{gil2002elaborar}, a pesquisa experimental consiste na definição de um objeto de estudo, na seleção de variáveis capazes de influenciá-lo e na determinação de formas de controle e observação dos efeitos produzidos por essas variáveis. Com base nessa definição, o presente trabalho caracteriza-se como uma pesquisa de natureza experimental, uma vez que modelos de tradução automática neural baseados na arquitetura Transformer serão comparados e submetidos a diferentes cenários e bases de dados.

Além disso, de acordo com \cite{da2002apostila}, a pesquisa quantitativa utiliza linguagem matemática para descrever fenômenos e analisar relações entre variáveis. Nesse contexto, este trabalho também possui caráter quantitativo, pois os modelos serão avaliados por meio de métricas como BLEU, BERTScore e chrF, bem como por medidas de eficiência computacional, obtidas a partir de experimentos controlados.


\section{Modelos}

Nesta seção são descritas as principais características dos modelos selecionados para comparação na tarefa de tradução automática, incluindo informações como tamanho, finalidade e forma de acesso.

A seguir, são apresentados os modelos utilizados neste trabalho, bem como suas respectivas configurações e locais para obtenção:

\begin{itemize}
    \item \textbf{MarianMT}: O modelo utilizado será o \textit{opus-mt-tc-big-en-pt}, com aproximadamente 0,2 bilhões de parâmetros. Trata-se de um modelo específico para tradução do inglês para o português. O seguinte \href{https://huggingface.co/Helsinki-NLP/opus-mt-tc-big-en-pt}{link} direciona para a plataforma \textit{Hugging Face}, onde o modelo pode ser obtido para uso local.

    \item \textbf{M2M-100}: O modelo utilizado será o \textit{m2m100\_418M}, com aproximadamente 0,4 bilhões de parâmetros. Diferentemente do MarianMT, este modelo não requer a seleção de um par específico de línguas, pois foi projetado para tradução direta entre múltiplos idiomas. O seguinte \href{https://huggingface.co/facebook/m2m100_418M}{link} redireciona para a plataforma \textit{Hugging Face}, onde o modelo pode ser obtido para uso local.

    \item \textbf{NLLB}: O modelo utilizado será o \textit{nllb-200-distilled-600M}, com aproximadamente 0,6 bilhões de parâmetros. Assim como o M2M-100, este modelo permite tradução direta entre múltiplas línguas, sem a necessidade de uma língua intermediária. O seguinte \href{https://huggingface.co/facebook/nllb-200-distilled-600M}{link} redireciona para a plataforma \textit{Hugging Face}, onde o modelo pode ser obtido para uso local.

\end{itemize}

\section{Tipos de texto}
\label{Text-type}
Nesta seção são definidos os tipos de texto utilizados na avaliação dos modelos de tradução automática.

Os textos foram organizados em quatro categorias, considerando o tamanho e o domínio do conteúdo, conforme descrito a seguir:

\begin{itemize}
    \item \textbf{Textos curtos}: textos com até 15 palavras de domínio conversacional e informal;
    
    \item \textbf{Textos longos}: textos com mais de 15 palavras de domínio conversacional e informal;

    \item \textbf{Textos específicos}: textos associados a domínios especializados, como a área de tecnologia da informação, também obtidos a partir de bases de dados públicas.

\end{itemize}


\section{Métricas de Avaliação}

Nesta seção tem como objetivo apresentar as ferramentas utilizadas para o cálculo das métricas BLEU, BERTScore e chrF e avaliação das traduções geradas pelos modelos.

\subsection{BLEU}

O cálculo da métrica BLEU, após a geração da tradução pelo modelo, será realizado utilizando a biblioteca da linguagem Python denominada NLTK (\textit{Natural Language Toolkit}). Essa biblioteca disponibiliza diversas funcionalidades voltadas ao processamento de linguagem natural, incluindo a função \textit{sentence\_bleu}, utilizada para calcular o valor da métrica BLEU a nível de sentença.

Essa função recebe como entrada uma ou mais sentenças de referência e a sentença gerada pelo modelo, retornando um valor numérico que representa o grau de similaridade entre elas com base em n-gramas.

Além disso, será utilizado um método de suavização (\textit{smoothing}) para evitar valores nulos em casos de baixa sobreposição entre os n-gramas da sentença gerada e da referência.

Abaixo é apresentado um exemplo de cálculo da métrica BLEU utilizando a função \textit{sentence\_bleu} da biblioteca NLTK:
\\
\\
\\
\\
\\
\\
\\
\\
\\
\\
\\
\begin{listing}
\begin{python}
from nltk.translate.bleu_score import sentence_bleu, SmoothingFunction

referencia = [['Uma', 'raposa', 'pula', 'sobre', 'o', 'cachorro', '.']]
geradoPeloModelo = ['Uma', 'pula','raposa', 'sobre', 'o', 'cachorro', ',']

smooth = SmoothingFunction().method1

print(sentence_bleu(referencia, geradoPeloModelo, smoothing_function=smooth))
\end{python}
\caption{Exemplo de cálculo da métrica BLEU com NLTK}
\label{code:bleu_example}
\end{listing}

\subsection{BERTScore}

O cálculo da métrica BERTScore será realizado utilizando a biblioteca da linguagem Python denominada \textit{bert-score}. Essa biblioteca permite calcular a similaridade entre sentenças por meio da comparação de \textit{embeddings} contextuais gerados a partir de modelos pré-treinados.

A função \textit{score} é responsável por receber como entrada a sentença gerada pelo modelo e a sentença de referência, retornando três medidas principais: precisão, revocação e F1-score.Para este trabalho, será utilizado o F1-score médio como medida final do BERTScore, por representar um equilíbrio entre precisão e revocação.

Abaixo é apresentado um exemplo de cálculo da métrica BERTScore utilizando a função \textit{score} da biblioteca \textit{bert-score}:

\begin{listing}
\begin{python}
from bert_score import score

geradoPeloModelo = ["O tempo esta congelante hoje."]
referencia = ["Hoje esta frio."]

PRECISAO, REVOCACAO, F1_SCORE = score(geradoPeloModelo, referencia, lang="pt", verbose=True)

print(F1_SCORE.mean())

\end{python}
\caption{Exemplo de cálculo da métrica BERTScore com a biblioteca bert-score}
\label{code:bleu_example}
\end{listing}

\subsection{chrF}

Para o cálculo da metrica chrF será realizada utilizada a biblioteca da linguagem python (\textit{sacrebleu}). Essa biblioteca possui funções que realizam o cálculo de métricas para avaliação de traduções automáticas, e dentre elas a função (\textit{sentence\_chrf}), cujo objetivo é realizar o cálculo da métrica chrF dado a sequência gerada pelo modelo e a sequência de referência.


A função (\textit{sentence\_chrf}) recebe dois argumentos, o primeiro sendo a sequência gerada pelo modelo, já o segundo argumento, a sequência de referência. Como resposta, a função retorna um valor entre 0 e 100, indicando o valor da métrica chrF

Abaixo é apresentado um exemplo de cálculo da métrica chrF utilizando a função \textit{sentence\_chrf} da biblioteca \textit{sacrebleu}:
\\
\begin{listing}
\begin{python}
import sacrebleu

geradoPeloModelo = "The cat is on the mat."
referencia = ["The cat is on the mat."]

chrf = sacrebleu.sentence_chrf(geradoPeloModelo, referencia)

print(chrf)

\end{python}
\caption{Exemplo de cálculo da métrica chrF com a biblioteca sacrebleu}
\label{code:bleu_example}
\end{listing}

\section{Ambiente Computacional}

Com o objetivo de garantir um ambiente controlado e uniforme para a comparação entre os modelos de tradução, optou-se pela utilização de um ambiente virtual em nuvem (\textit{cloud}) para execução dos experimentos.A plataforma escolhida foi a \textbf{DigitalOcean}, devido à disponibilidade de instâncias com cobrança por hora e oferta de máquinas com alto desempenho, o que possibilita maior flexibilidade na execução dos testes.

A Tabela abaixo apresenta as especificações da máquina virtual utilizada nos experimentos:

\begin{table}[H]
\centering
\begin{tabular}{ |p{3cm}|p{3cm}|p{3cm}|p{3cm}| }
\hline
\multicolumn{1}{|p{3cm}|}{\centering\textbf{Memória RAM}} & \multicolumn{1}{p{3cm}|}{\centering\textbf{Núcleos (vCPU)}} & \multicolumn{1}{p{3cm}|}{\centering\textbf{Transferência de dados}} & \multicolumn{1}{p{3cm}|}{\centering\textbf{Armazenamento}} \\
\hline
\multicolumn{1}{|p{3cm}|}{\centering 16 GB} & \multicolumn{1}{p{3cm}|}{\centering 8 vCPUs} & \multicolumn{1}{p{3cm}|}{\centering 6 TB} & \multicolumn{1}{p{3cm}|}{\centering 100 GB} \\
\hline
\end{tabular}
\caption{Especificações da máquina virtual utilizada}
\label{tab:ambiente_computacional}
\end{table}

Cabe destacar que, por se tratar de um ambiente virtualizado, informações detalhadas sobre o hardware físico, como modelo do processador ou tipo de memória, não são explicitamente fornecidas pela plataforma. Isso ocorre porque a instância virtual pode ser alocada dinamicamente sobre diferentes máquinas físicas, dependendo da disponibilidade no momento da criação.

\section{Desempenho Computacional}

Nesta seção são descritas as métricas e ferramentas utilizadas para avaliação da eficiência computacional dos modelos de tradução, considerando principalmente o uso de CPU, tempo de execução e consumo de memória RAM.

\subsection{Uso de CPU e Tempo de Execução}

Neste experimento, serão considerados dois indicadores principais relacionados à CPU: o tempo médio de execução e o uso médio de CPU durante a tarefa de tradução.

Para a medição do tempo de execução, será utilizada a ferramenta de linha de comando \textit{hyperfine}. Essa ferramenta é amplamente utilizada para a realização de benchmarks, permitindo medir o tempo total, média e desvio padrão de execução de um programa ao longo de múltiplas execuções.

Abaixo é apresentado um exemplo de uso da ferramenta \textit{hyperfine}:

\begin{listing}
\begin{python}
hyperfine 'python3 script.py'
\end{python}
\caption{Exemplo de uso da ferramenta \textit{hyperfine}}
\label{code:hyperfine}
\end{listing}

Para a medição do uso médio de CPU durante a execução dos modelos, será utilizada a ferramenta \textit{pidstat}, nativa de sistemas operacionais Linux. Essa ferramenta permite coletar estatísticas detalhadas sobre o uso de recursos por processos individuais.

A seguir, um exemplo de utilização da ferramenta \textit{pidstat}:

\begin{listing}
\begin{python}
pidstat -u -e python3 test.py
\end{python}
\caption{Exemplo de uso da ferramenta \textit{pidstat}}
\label{code:pidstat}
\end{listing}

As principais flags utilizadas são:

\begin{itemize}
    \item \textbf{-u}: exibe estatísticas de utilização da CPU;
    \item \textbf{-e}: executa o comando especificado e monitora seu consumo de recursos.
\end{itemize}

\subsection{Consumo de Memória RAM}

Para a análise de memória RAM, serão considerados dois indicadores: o uso médio e o uso máximo (pico) durante a execução da tarefa de tradução. O uso médio de memória será obtido por meio da ferramenta \textit{pidstat}, utilizando a flag \textit{-r}, que permite monitorar o consumo de memória ao longo da execução do processo.

Já o uso máximo de memória (pico) será medido utilizando a ferramenta \textit{/usr/bin/time}, também disponível em sistemas Linux. Essa ferramenta fornece informações detalhadas sobre a execução de programas, incluindo o consumo máximo de memória.

Abaixo é apresentado um exemplo de uso da ferramenta \textit{time}:

\begin{listing}
\begin{python}
/usr/bin/time -v python3 script.py
\end{python}
\caption{Exemplo de uso da ferramenta \textit{time}}
\label{code:time}
\end{listing}

A flag \textit{-v} permite a exibição de informações detalhadas, incluindo o valor de memória máxima utilizada durante a execução do programa.

\section{Base de Dados}
\label{Database}
Esta seção apresenta as bases de dados utilizadas nos experimentos de avaliação dos modelos de tradução automática.

\subsection{Textos Conversacionais e Informais}

Para a composição da base de textos conversacionais e informais, será utilizado o repositório OPUS (\textit{Open Parallel Corpus}), mais especificamente o conjunto de dados \textit{OpenSubtitles v2024}, contendo pares de tradução entre inglês e português brasileiro.

Essa base é composta majoritariamente por diálogos extraídos de legendas de filmes e séries, sendo adequada para representar linguagem informal e conversacional.

A base de dados pode ser acessada por meio do seguinte \href{https://opus.nlpl.eu/datasets/OpenSubtitles?pair=en&pt-BR}{link}.

A Tabela \ref{tab:OpenSubtitlesv2024} apresenta exemplos de pares de frases presentes na base de dados.

\begin{table}[H]
\centering
\begin{tabular}{ |>{\centering\arraybackslash}p{6cm}|>{\centering\arraybackslash}p{6cm}| }
\hline
\textbf{Frase em inglês} & \textbf{Tradução para português brasileiro} \\
\hline
Did you do this, Dr. Ivins? & Você fez isso, Dr. Ivins? \\
\hline
Hello, mother. Hello, Bruce. & Olá, mãe. Olá, Bruce. \\
\hline
Hello, papa. What are you doing? & Olá, papai. O que você está fazendo? \\
\hline
\end{tabular}
\caption{Exemplo de pares de tradução da base OpenSubtitles v2024}
\label{tab:OpenSubtitlesv2024}
\end{table}

\subsection{Textos Específicos}

Para a avaliação em textos de domínio específico, serão utilizadas bases de dados do repositório OPUS com conteúdos técnicos. Foram selecionados os conjuntos \textit{MDN Web Docs}, referente à documentação de tecnologias web mantida pela organização Mozilla, e \textit{PHP}, contendo documentação da linguagem de programação PHP.

As bases podem ser acessadas por meio do seguinte \href{https://opus.nlpl.eu/corpora-search/en&pt-BR}{link}.

A Tabela \ref{tab:MDNwebdocs} apresenta exemplos de pares de frases da base \textit{MDN Web Docs}:

\begin{table}[H]
\centering
\begin{tabular}{  |>{\centering\arraybackslash}p{6cm}|>{\centering\arraybackslash}p{6cm}| }
\hline
\textbf{Frase em inglês} & \textbf{Tradução para português brasileiro} \\
\hline
Present, accept, interpret, calculate, repeat & Apresentar, aceitar, interpretar, calcular, repetir \\
\hline
The specifics depend on the game. & As especificidades dependem do jogo. \\
\hline
Some games drive this cycle by user input. & Alguns jogos guiam este ciclo pela entrada do usuário. \\
\hline
\end{tabular}
\caption{Exemplo de pares de tradução da base MDN Web Docs}
\label{tab:MDNwebdocs}
\end{table}

A Tabela \ref{tab:PHP} apresenta exemplos de pares de frases da base \textit{PHP}:

\begin{table}[H]
\centering
\begin{tabular}{  |>{\centering\arraybackslash}p{6cm}|>{\centering\arraybackslash}p{6cm}| }
\hline
\textbf{Frase em inglês} & \textbf{Tradução para português brasileiro} \\
\hline
The documentation is stored in the phpdoc module. & A documentação é armazenada no módulo phpdoc. \\
\hline
How to help improve the documentation & Como ajudar a melhorar a documentação \\
\hline
Classify the bug as "Documentation Problem". & Classifique o problema como "Documentation Problem". \\
\hline
\end{tabular}
\caption{Exemplo de pares de tradução da base PHP}
\label{tab:PHP}
\end{table}


\section{Procedimento de Execução}

\begin{figure}[H]
    \centering
    \includegraphics[height=0.7\textwidth]{Imagens/process.pdf}    \caption{\label{fig:procedimento} Procedimento}
    \legend{Fonte: autor}
\end{figure}

\subsection{Seleção da base de dados}

Conforme descrito na Seção \ref{Database}, a base de dados utilizada neste trabalho será a OPUS (Open Parallel Corpus), escolhida por disponibilizar corpora paralelos públicos para diferentes pares de idiomas. A partir dessa base serão selecionados os conjuntos de textos utilizados nos experimentos, incluindo textos curtos, textos longos e textos de domínio específico.

\subsection{Seleção dos textos}

Conforme apresentado na Seção \ref{Text-type}, os experimentos considerarão três categorias de textos: textos curtos, textos longos e textos de domínio específico. Para cada categoria, serão selecionados pares de sentenças em inglês e português brasileiro.

Durante o processo de seleção, serão removidos dados que apresentem inconsistências que possam comprometer a avaliação dos modelos, tais como traduções incorretas, caracteres inválidos, símbolos não pertencentes ao idioma de destino, excesso de espaços em branco e quebras de linha desnecessárias.

Pretende-se utilizar aproximadamente 100 pares de sentenças para cada categoria de texto, com o objetivo de criar uma diversidade entre os textos escolhidos para tradução. A seleção, filtragem e categorização definitiva dos dados serão realizadas durante a etapa experimental do trabalho.

\subsection{Tradução gerada e métricas}

Após a seleção e limpeza dos dados, cada modelo de tradução irá executar na mesma máquina virtual, sob as mesmas configurações de hardware e receberá exatamente as mesmas sentenças de entrada em inglês, garantindo condições equivalentes durante os experimentos. As traduções geradas serão então comparadas com suas respectivas traduções de referência em português brasileiro.

Durante o processo de inferência, serão coletadas métricas relacionadas ao desempenho computacional, incluindo utilização de CPU, consumo de memória RAM e tempo de execução. Posteriormente, as traduções geradas serão avaliadas por meio das métricas BLEU, BERTScore e chrF, com o objetivo de analisar a qualidade das traduções produzidas por cada modelo.

Para reduzir possíveis variações causadas pelo ambiente de execução, cada experimento será executado três vezes para o mesmo conjunto de sentenças. Ao final, será calculado o valor médio das métricas obtidas em cada execução, bem como a variação observada entre os resultados, permitindo uma análise mais consistente do desempenho dos modelos avaliados.


\section{Análise dos dados}

Esta seção tem como objetivo descrever os procedimentos adotados para a análise dos dados coletados a partir dos experimentos de tradução automática. A análise considera tanto a qualidade das traduções geradas quanto o desempenho computacional dos modelos avaliados.

\subsection{Comparação entre modelos}

A comparação entre os modelos de tradução será realizada a partir da utilização de um mesmo conjunto de entradas, composto por pares de sentenças paralelas, sendo o inglês a língua de origem e o português brasileiro a língua de destino. Cada modelo receberá exatamente as mesmas sentenças de entrada, a fim de garantir a comparabilidade dos resultados.

Para cada sentença de entrada, será gerada uma tradução por cada modelo avaliado. As saídas produzidas serão então comparadas com as respectivas traduções de referência por meio das métricas BLEU, BERTScore e chrF. Além disso, serão coletadas métricas relacionadas ao desempenho computacional, incluindo uso de memória RAM, utilização de CPU e tempo de execução durante o processo de inferência.

\subsection{Comparação entre tipos de texto}

A análise também considera o impacto das características dos textos de entrada no desempenho dos modelos. Para isso, os dados serão organizados em duas dimensões principais: tamanho do texto e domínio do conteúdo.

Em relação ao tamanho, as sentenças serão classificadas em categorias previamente definidas, como textos curtos e textos longos, com base no número de palavras ou tokens. Já em relação ao domínio, os textos serão agrupados em categorias distintas, como textos conversacionais e textos técnicos, sendo estes últimos compostos por conteúdos relacionados à linguagem de programação PHP e à documentação web da fundação Mozilla.

Para cada categoria, um conjunto equivalente de sentenças será selecionado e submetido a todos os modelos avaliados. As métricas de qualidade (BLEU, BERTScore e chrF) e de desempenho computacional serão calculadas individualmente para cada categoria.

Os resultados serão analisados de forma comparativa, por meio de tabelas e representações gráficas, permitindo observar o comportamento dos modelos em diferentes cenários.

\subsection{Correlação entre métricas}

Além da análise individual das métricas, este trabalho também investiga a relação entre elas. As métricas BLEU, BERTScore e chrF possuem fundamentos distintos: enquanto o BLEU é baseado na sobreposição de n-gramas, o BERTScore considera similaridade semântica por meio de representações vetoriais, e o chrF avalia correspondências em nível de caracteres. Dessa forma, é esperado que cada métrica capture diferentes aspectos da qualidade da tradução.

Para analisar a relação entre essas métricas, será realizada uma análise de correlação estatística entre seus valores, considerando os resultados obtidos para todas as sentenças e modelos avaliados. Essa análise tem como objetivo identificar o grau de concordância entre as métricas, bem como possíveis divergências em suas avaliações.




